% 三种距离：p=(0,0), q=(4,0), r=(3,3)。d1(p,r)=6（折线）, d2(p,r)=3sqrt2≈4.243（直线）, d_inf(p,r)=3；三种距离 p 到 q 均为 4。
\begin{tikzpicture}[scale=.75,
  pt/.style={circle,fill=FigureStroke,inner sep=0pt,minimum size=4pt},
  note/.style={font=\FigureNoteFont\fontsize{7.5}{9}\selectfont,align=center}]
  \draw[FigureGrid,line width=.4pt] (-.3,-.3) grid (4.6,3.6);
  \draw[->,FigureStroke,line width=.6pt] (-.3,0)--(4.8,0) node[right]{$x_1$};
  \draw[->,FigureStroke,line width=.6pt] (0,-.3)--(0,3.8) node[above]{$x_2$};
  \draw[FigureGreen,line width=1.1pt] (0,0)--(3,0)--(3,3);
  \draw[FigureBlue,line width=1.1pt] (0,0)--(3,3);
  \draw[FigureYellow,line width=2pt,opacity=.9] (3,0.04)--(3,2.96);
  \node[pt,label=below left:$\bm p$] at (0,0){};
  \node[pt,label=below:$\bm q$] at (4,0){};
  \node[pt,label=above right:$\bm r$] at (3,3){};
  \node[note,text=FigureBlue,anchor=east] at (1.3,1.9) {$d_2=3\sqrt2$};
  \node[note,text=FigureGreen,anchor=north] at (1.5,-.05) {$d_1=3+3=6$};
  \node[note,text=FigureYellow!70!FigureInk,anchor=west] at (3.1,1.5) {$d_\infty=3$};
  \node[note,text=FigureMuted,anchor=west] at (3.6,.45) {$\bm p$到$\bm q$：都为$4$};
\end{tikzpicture}
